\documentclass[../../../include/open-logic-section]{subfiles}

\begin{document}

\olfileid{sth}{choice}{banach}
\olsection{The Banach--Tarski（巴拿赫--塔斯基）悖论}

正如 Boolos（布洛斯）尝试为替换公理辩护一样，我们或许也可以尝试诉诸\emph{外在理由}来为选择公理辩护（参见 \olref[replacement][extrinsic]{sec}）。毕竟，采用选择公理有许多理想的推论：能够比较每个基数；能够良序化每个集合；能够将基数视为特定种类的序数；等等。

然而，有时人们声称选择公理有\emph{不良的}后果。这主要源于\cite{BanachTarski1924}的一个结果。

\begin{thm}[Banach--Tarski 悖论（在 $\ZFC$ 中）]
任何球体都可以被分解为有限个部分，这些部分可以通过（旋转和平移）重新组装，形成该球体的两个副本。
\end{thm}

\noindent
乍看之下，这有些令人惊异。显然，两个球体的体积是原球体的\emph{两倍}。但是刚体运动——旋转和平移——不改变体积。因此，Banach--Tarski 悖论看起来允许我们凭空变出新的物质。

更夸张的是，\footnote{参见 \citet[Theorem 3.12]{Wagon2016}。}类似的推理表明，一颗豌豆可以被切成有限个部分，然后这些部分可以（通过旋转和平移）重新组装，形成一个形状和大小如同大本钟的物体。

然而，仅凭 $\ZF$ 自身，这些都不成立。\footnote{不过，Banach--Tarski 悖论可以用严格弱于选择公理的原则证明；参见 \citet[303]{Wagon2016}。}因此我们面临一个抉择：拒绝选择公理，或者学会与这个“悖论”共存。

我们的建议是，应当学会与这个“悖论”共存。事实上，我们根本不认为它算得上什么悖论。特别是，我们看不出它与以下任何结果相比更为悖谬：\footnote{\citet[276--7]{Potter2004}, \citet[16]{Weston2003}, \citet[31, 308--9]{Wagon2016} 使用其他例子提出了类似观点。}

\begin{enumerate}
	\item 区间 $(0,1)$ 中的点和 $\Real$ 中的点一样多。
	\\\emph{证明}: 考虑 $\tan(\pi(r-\nicefrac{1}{2})))$。
	\item 一条线上的点和一个正方形中的点一样多。
	\\参见 \olref[his][set][pathology]{sec} 和 \olref[his][set][cantorplane]{sec}。
	\item 存在空间填充曲线。
	\\参见 \olref[his][set][pathology]{sec} 和 \olref[his][set][hilbertcurve]{sec}。
\end{enumerate}

这三个结果都不需要选择公理。事实上，我们现在只是把它们视为令人惊讶、奇妙的数学片段。也许我们应该对 Banach--Tarski 悖论采取同样的态度。

当然，这里需要一个技术性观察；但这只需保持清醒。刚体运动保持体积。因此，球体被分解成的那五个部分\footnote{我们陈述悖论时用了“有限个部分”。实际上，\citet{Robinson1947}（罗宾逊）证明分解可以用\emph{五}个部分完成（但不能再少）。证明参见 \citet[pp.~66--7]{Wagon2016}。}不可能都是\emph{可测的}。粗略地说，为这些单独的部分赋予体积是没有意义的。你应该把它们想象成无法描绘的、点的“无限散布”。如今，构想出这样的“无限散布”集合也许是“怪异”的。但它们的存在性似乎源于\stagesacc{}所体现的原则：你应该构造出所有可能的、由先前已得集合组成的汇集。

如果所有这些都不能说服你，这里还有一个支持接纳 Banach--Tarski 悖论的最终（外在）理由。它直接引出了有史以来最好的数学笑话：

\begin{enumerate}
	\item[] \emph{问}。“Banach--Tarski”的一个字母重排词是什么？
	\item[] \emph{答}。“Banach--Tarski Banach--Tarski”。
\end{enumerate}

\end{document}